\section{The trimming}
\label{sec:trimming}

Once we know that at least a fraction $\widehat\alpha$ of each population has no counterpart in the other, the next
question is which observations these are. By Proposition~\ref{prop:bangbang} the minimizer of \eqref{eq:qp} answers it
with a pair of trimmed samples, which keep at full weight the observations on one side of a cut of the witness and split
weight among those at the cut. They converge, at a fixed level to the best trimming at that level and at the estimated
level to the common core. The rate we prove at a fixed level assumes a witness with a density at the cut, condition
\eqref{eq:qdens} below, as for the Gaussian kernel (Remark~\ref{rem:qdens} of the supplement). With the Laplace kernel
on the line the population witness equals its cut on a set of $P$-mass at least $\alpha$ (Section~\ref{sec:excessproofs}
of the supplement), so \eqref{eq:qdens} fails at every level, Theorem~\ref{cor:trimrate} does not apply, and the
trimming is decided among tied observations; Section~\ref{sec:trimest}, in MMD and without a rate, covers every kernel
of Assumption~\ref{ass:regularity}.

\subsection{A parametric rate at a fixed level}

Fix $\alpha$ with $D=D_\alpha(P,Q)>0$, which under (A3) means $\alpha<\alpha^\star$, let $d^\star$ be the population
witness of Section~\ref{sec:statistic}, with direction $f^\star=d^\star/D$, and let $t_\alpha$ be a $(1-\alpha)$-quantile
of $d^\star(X)$ under $P$. The population version of Proposition~\ref{prop:bangbang} holds: every optimal trimming
density of $P$ has the threshold form (Lemma~\ref{lem:popthreshold} of the supplement)
\begin{equation}
g^\star=\beta\,\mathbf 1\{d^\star<t_\alpha\}+\vartheta\,\mathbf 1\{d^\star=t_\alpha\},
\qquad 0\le\vartheta\le\beta\ \text{a function with}\ \E_P[g^\star]=1 .
\label{eq:popdensity}
\end{equation}
The observations that an optimal trimming keeps at full weight form the \emph{kept sets}
\[
A^\star\coloneqq\{x:\ d^\star(x)<t_\alpha\}
\qquad\text{and}\qquad
\widehat A_n\coloneqq\{x:\ \widehat d(x)<\widehat t_n\},
\]
population and empirical, whose symmetric difference $\widehat A_n\triangle A^\star$ measures the error, and we let
$\widehat g_n$ be the density \eqref{eq:popdensity} with $(\widehat d,\widehat t_n)$ in place of $(d^\star,t_\alpha)$
and any tie function $0\le\widehat\vartheta\le\beta$, a rule that need not integrate to one under $P$. Level sets are
stable where the witness has a density at the cut, which we ask at both cuts:
\begin{gather}
\tag{Q}\label{eq:qdens}
\parbox[t]{0.84\textwidth}{\raggedright the law of $d^\star(X)$ under $P$ has, on $[t_\alpha-r_Q,t_\alpha+r_Q]$ for some
$r_Q>0$, a density with values in $[c_f,C_f]$, where $0<c_f\le C_f<\infty$;}\\[4pt]
\tag{Q$'$}\label{eq:qdensY}
\parbox[t]{0.84\textwidth}{\raggedright the law of $d^\star(Y)$ under $Q$ has, on $[u_\alpha-r_Q,u_\alpha+r_Q]$, a
density at least $c'_f>0$, where $u_\alpha$ is an $\alpha$-quantile of $d^\star(Y)$.}
\end{gather}
Condition \eqref{eq:qdens} is the classical quantile-density condition for the cut of the $X$ sample, and
\eqref{eq:qdensY} its lower half for the cut of the $Y$ sample. Laws with finite support, such as those of
Example~\ref{ex:separated} with finitely supported $R$, $P_0$, $Q_0$, violate \eqref{eq:qdens}, so the margin does not imply the density conditions. The rate comes from a localization argument of the kind behind rates for M-estimators
\citep[Theorem~3.2.5]{vandervaart1996weak}: the separation of Section~\ref{sec:ramp} falls quadratically away from its
maximizer $f^\star$, while its fluctuation grows linearly in the distance to $f^\star$, times $(n\wedge m)^{-1/2}$.

\begin{theorem}[A parametric rate for the witness and the kept set]
\label{cor:trimrate}
Let (A1) and (A2) of Assumption~\ref{ass:regularity} hold, fix $\alpha\in(0,1)$ with $D=D_\alpha(P,Q)>0$, and suppose
\eqref{eq:qdens} and \eqref{eq:qdensY}. Put $V_\alpha=\kappa\beta+[\beta/\min(c_f,c'_f)]^{1/2}$ and
$W_\alpha=C_f(\kappa V_\alpha+c_f^{-1})$. Then, with $(P,Q,k,\alpha)$ fixed and $n\wedge m\to\infty$,
\begin{gather*}
\big\|\widehat f-f^\star\big\|_\Hh=O_\Pp\Big(\frac{V_\alpha}{D\sqrt{n\wedge m}}\Big),
\qquad
\|\widehat d-d^\star\|_\Hh=O_\Pp\Big(\frac{V_\alpha}{\sqrt{n\wedge m}}\Big),\\
P(\widehat A_n\triangle A^\star)=O_\Pp\Big(\frac{W_\alpha}{\sqrt{n\wedge m}}\Big),
\qquad
\frac1\beta\int|\widehat g_n-g^\star|\,dP=O_\Pp\Big(\frac{W_\alpha}{\sqrt{n\wedge m}}\Big),
\end{gather*}
with constants in the $O_\Pp$ that depend only on the probability level.
\end{theorem}

The rate is that of a sample quantile, free of the dimension of $\mathcal X$, because the witness converges in $\Hh$,
as a mean of feature vectors, and the kept set inherits its rate through the density at the cut. As
$\alpha\uparrow\alpha^\star$ the bound on the direction grows like $1/D$, and the sample size from which the bounds hold,
not explicit (Section~\ref{app:trimming} of the supplement states what it depends on), grows as $D$ vanishes. The kept set then has in general no limit, because the common
core keeps only a fraction of $P$ on the set $\{0<q<p\}$, and an indicator of a set cannot converge to such a
fractional weight. With the Wasserstein distance, optimal trimmings are consistent and computable
\citep{alvarez2011incomplete,agullo2018trimming}, and rates are known only toward a known target
\citep{delbarrio2013rates}.

\subsection{The trimming at the estimated level}
\label{sec:trimest}

In practice the trimming is read at a data-driven level, the lower bound $\widehat\alpha$, where
Theorem~\ref{cor:trimrate} does not apply: its level is fixed in advance, and its bounds degrade as $D_\alpha(P,Q)$
vanishes near $\alpha^\star$. There the trimmed sample still converges as a measure. The common core $R^\star$ is the unique common
trimming at $\alpha^\star$, and
\[
\Omega(\epsilon)=\sup\big\{\|a-\mu_{R^\star}\|_\Hh:\ a\in\mathcal C_{\alpha^\star}(P),\
\dist\big(a,\mathcal C_{\alpha^\star}(Q)\big)\le\epsilon\big\}
\]
is a modulus of the contact of the two regions at $\alpha^\star$.

\begin{proposition}[The trimmed sample at the data-driven level]
\label{prop:alphahat}
Let Assumption~\ref{ass:regularity} hold and let $0<\alpha^\star<1$. Then $\Omega(\epsilon)\to0$ as
$\epsilon\downarrow0$. Fix $\delta'\in(0,1)$ and put $r_n=\kappa(1-\alpha^\star)^{-1}[4n^{-1/2}+(2\log(2/\delta')/n)^{1/2}]$,
and $r_m$ likewise. There is an event $E$ of probability at least $1-2\delta'$, not depending on the level, on which for
every level $\tilde\alpha\le\alpha^\star$, possibly chosen from the data, and every $(w,v)\in\mathcal T_{\tilde\alpha}$,
\[
\MMD\big(P_n^w,R^\star\big)\le r_n+\Omega\big(\widehat D_{\tilde\alpha}+r_n+r_m\big).
\]
The same holds for the second sample, with $P$ and $Q$ exchanged.
\end{proposition}

On $E$ each empirical region is within $r_n$, respectively $r_m$, of its population counterpart, at all levels up to
$\alpha^\star$ at once, by \eqref{eq:hausdorff}. The bound covers the trimming read at $\widehat\alpha$ or at any other
valid lower bound. It carries no rate unless $\Omega$ is known, and it places the reweighted sample near the common core
without saying which observations are kept.

\begin{corollary}[The trimmed sample at the lower bound]
\label{cor:alphalo}
Let Assumption~\ref{ass:regularity} hold, let $0<\alpha^\star<1$, and compute $\widehat\alpha$ at level
$\delta$ as in Proposition~\ref{prop:loconsistency}, on grids fixed in advance whose mesh tends to zero, with
$\log J=o(n\wedge m)$. Then for every $\epsilon>0$ and $\delta'\in(0,1)$, for all $n,m$ large enough, with probability
at least $1-\delta-\delta'$ every $(w,v)\in\mathcal T_{\widehat\alpha}$ has
$\MMD(P_n^w,R^\star)\le\epsilon$ and $\MMD(Q_m^v,R^\star)\le\epsilon$.
\end{corollary}

By Lemma~\ref{lem:pathconvex}(ii) of the supplement the path falls to zero at most linearly, so a consistent
$\widehat\alpha$ has a small trimmed distance whatever the laws and no margin is needed. The proof uses
$\widehat\alpha$ only through its validity and its consistency, so the level may come from a bound computed with
another kernel than the one that trims, as in Section~\ref{sec:appcohort}.
